\chapter{Actual renormalization: effective actions and scale}
\label{rg:chapter}
Renormalization changes the description while preserving a specified set of
physical predictions. To make that claim testable, we carry the effective
action, its constant, observables and scale convention. Neural state
compression and principal-axis alignment can be useful auxiliary techniques,
but their training losses do not replace these objects.
% Generated shared actual RG material.
\section{Actual renormalization retains the effective action}
\label{sec:actual-rg}
A compression score does not specify an RG transformation. Following the
real-space formulation of Wilson and Kogut \cite{WilsonKogutRG}, we specify
the eliminated variables, effective action, normalization and scale map.
Our exact finite decimation is distinct from an approximate tensor-network
truncation such as Levin--Nave TRG \cite{LevinNaveRG}.

\begin{definition}[Finite classical RG step]
Let the fine configuration be $(s,t)\in S\times E $, with finite nonempty
sets $ S,E $, and dimensionless log Boltzmann weight $ V(s,t)=-\beta H(s,t)$.
Define
\[
 (\mathcal RV)(s)=\log\sum_{t\in E}e^{V(s,t)},\qquad
 Z(V)=\sum_{s,t}e^{V(s,t)}.
\]
The full action includes its additive constant. A truncation or a learned
approximation is a further operation with its own error, not this identity.
\end{definition}
\begin{proposition}[Partition, observables and composition]\label{rg:exact}
One has $ Z(\mathcal RV)=Z(V)$, and the normalized expectation of every
observable depending only on $ s $ agrees in the two descriptions. Eliminating
$ E_1$ and then $ E_2$ equals eliminating $ E_1\times E_2$ at once, with the
same tensor labels and inverse temperature.
\end{proposition}
\begin{proof}
Exponentiating the definition gives $ e^{\mathcal RV(s)}=\sum_te^{V(s,t)}$.
Sum over $ s $, optionally multiplying by an observable $ O(s)$, and divide
both sides by the same strictly positive $ Z $. Reordering finite sums proves
the composition statement. No projection onto pair interactions is used.
\end{proof}

For an even homogeneous periodic Ising chain, write the transfer matrix as
$ T_{ab}=\exp(C+Kab+h(a+b)/2)$, $ a,b\in\{-1,1\}$.
Its partition function is $\operatorname{tr}T^N $. Summing alternate spins
replaces $ T $ by $ T^2$ and $ N $ by $ N/2$; the physical spacing doubles.
Writing $ P=T^2$, its positive entries determine
\[
 K'=\tfrac 14\log\frac{P_{++}P_{--}}{P_{+-}^2},\quad
 h'=\tfrac 12\log\frac{P_{++}}{P_{--}},\quad
 C'=\tfrac 14\log(P_{++}P_{--}P_{+-}^2).
\]
These three parameters reproduce every entry of $ P $, so no normalization
is lost. At $ h=0$, $\tanh K'=\tanh^2 K $ and the eigenvalue-ratio correlation length of the homogeneous extension
in blocked units halves. Its product with the spacing stays fixed. The
fixed points in $ t=\tanh K $ are $ t=0,1$; the latter is zero temperature,
not a positive-temperature transition. Odd dessin examples below use the
full-action definition, not this even-chain shortcut.

\subsection{Genuine dessin incidence models}
For odd $ m $, the clean dessin $ y^2=x^m-1$, $\beta=x^m $ has two black
vertices, $ m $ white vertices and one face; its genus is $(m-1)/2$.
Its incidence graph is $ K_{2,m}$ with $2 m $ clean edges. Its infinite
locally gentle order quiver instead has $ m $ vertices and completed
DVR-rank $2 m^2$. We place physical spins $ B_+,B_-,s_0,\ldots,s_{m-1}$
on the incidence vertices; this is a specified finite spin model, not a
faithful representation of that infinite algebra or order. The abstract
incidence graph alone does not recover the ribbon embedding or genus.

With $\eta_j=(-1)^j $ and fixed nonuniform fields $ h_j $, choose
\[
 V=J_+B_+\sum_js_j+J_-B_-\sum_j\eta_js_j
       +bB_+-\tfrac 35 bB_-+\sum_jh_js_j.
\]
Summing the two black spins gives the \emph{entire} effective action
\[
 V'(s)=\sum_jh_js_j+
 \log\!\left[2\cosh\!\left(b+J_+\sum_js_j\right)\right]+
 \log\!\left[2\cosh\!\left(-\tfrac 35 b+J_-\sum_j\eta_js_j\right)\right].
\]
Every one of the $2 m $ physical edges is present, although its couplings
belong to a tied three-parameter family. The Walsh expansion
$ V'(s)=\sum_{A\subseteq\{0,\ldots,m-1\}}c_A\prod_{j\in A}s_j $,
with $ c_A=2^{-m}\sum_sV'(s)\prod_{j\in A}s_j $, is unique by Boolean
orthogonality. It retains the constant $ c_\varnothing $ and genuinely
generated higher-body interactions. Our genus-two and genus-four examples
use $ m=5,9$ and continue by exact nested marginalization. Dropping all terms
of degree greater than two is displayed as a measured approximation.

\subsection{Quantum thermal effective Hamiltonians}
For a finite Hermitian Hamiltonian $ H $ on nonzero finite-dimensional
$\mathcal H_S\otimes\mathcal H_E $
and $\beta>0$, define
\[
 W'=\operatorname{Tr}_E e^{-\beta H},\qquad
 H_{\mathrm{eff}}(\beta)=-\beta^{-1}\log W'.
\]
The matrix $ W'$ is strictly positive: for $ v\ne 0$, sum the strictly
positive quadratic forms of $ v\otimes e_j $ over an orthonormal basis of
$\mathcal H_E $. Its Hermitian spectral logarithm is therefore defined.
The identity $ e^{-\beta H_{\mathrm{eff}}}=W'$ proves partition and retained
\emph{equal-time} observable preservation by the partial-trace property.
Repeated traces compose. The full Pauli expansion, including its identity
coefficient, is retained. It is generally nonlocal and depends on $\beta $;
this statement does not preserve all fine real-time dynamics or supply a
local zero-temperature/thermodynamic RG fixed point.

\texttt{FiniteRG.lean} checks nine actual finite scalar marginal,
exponential-action, partition, observable, composition, error-sandwich and
transfer-trace statements. The quantum spectral-log and complete
thermodynamic/field-theory interfaces are not formalized by these lemmas.
These physics constructions give no arithmetic determinant comparison or
new GRH conclusion.

\begin{center}
\begin{tikzpicture}[scale=.85,every node/.style={font=\small}]
\node[draw,circle] (p) at (0,1.5) {$ B_+$};
\node[draw,circle] (n) at (0,-1.5) {$ B_-$};
\foreach \j in {0,...,4}{
 \node[draw,circle,fill=white] (w\j) at (2,{2-\j}) {$ s_{\j}$};
 \draw[brown] (p)--(w\j);\draw[blue] (n)--(w\j);
}
\draw[-{Latex},thick] (3,0)--(4.5,0);
\node[draw,rounded corners,text width=55 mm,align=center] at (8,0)
 {Sum both black spins\\$ V'=c_\varnothing+\sum_Ac_A\prod_{j\in A}s_j $\\All generated interactions retained};
\end{tikzpicture}
\end{center}
This incidence schematic displays the physical elimination, not a planar
uniformization of the genus-two surface. The explicit $\sigma,\alpha $ and
face cycle determine the ribbon embedding; the interactive laboratory
shows actual radial lifts and labels its four-real-coordinate projection.
% Generated shared actual RG material.
\section{Controlled errors and a trained full-action RG approximation}
\label{sec:rg-details}
\begin{proposition}[Classical action-error propagation]
If $|\widehat V(s,t)-V(s,t)|\le\varepsilon $ at every finite configuration,
then $|\mathcal R\widehat V(s)-\mathcal RV(s)|\le\varepsilon $.
The partition ratio lies in $[e^{-\varepsilon},e^{\varepsilon}]$.
Every further \emph{exact} marginalization preserves the same action bound.
\end{proposition}
\begin{proof}
The pointwise inequalities
$ e^{-\varepsilon}e^V\le e^{\widehat V}\le e^\varepsilon e^V $
survive finite summation. Taking logarithms proves the first assertion;
summing also over the retained variables proves the second. Iterate the
same argument. For normalized measures the likelihood ratio lies between
$ e^{-2\varepsilon}$ and $ e^{2\varepsilon}$, so a bounded observable error
is at most $\|O\|_\infty\min(2,e^{2\varepsilon}-1)$. Per-step learned
errors would add; a single initial error does not grow under subsequent
exact elimination. This scalar proof is not an assertion that the matrix
exponential is operator-monotone.
\end{proof}

\subsection{Executed finite calculations}
The implementation checks full enumeration against the closed black-spin
formula, Walsh reconstruction, raw partition equality, retained-spin
magnetizations and two-stage versus direct elimination. At the quantum
$\Gamma=0$, $\beta=1$ limit, every Pauli coefficient agrees with the
classical effective action, with the required minus sign between energy
and logweight. At nonzero transverse field the model is
\[
 H=-J_+Z_+\sum_jZ_j-J_-Z_-\sum_j\eta_jZ_j-bZ_+
       +\tfrac 35 bZ_--\sum_jh_jZ_j-\Gamma\sum_iX_i.
\]
At general $\beta $, the classical zero-transverse-field limit uses $\beta V $, including the fixed white fields, rather than scaling only the three varied controls.
For $ m=5$, its full dimension is 128 and the first retained dimension is 32.
Every one of the 1024 five-spin Pauli coefficients is computed; later
retained dimensions are 8,4,2. Spectral-log/exponential and full Pauli
reconstruction residuals, exact-trace composition and retained-observable
comparisons are recorded separately. Floating-point residuals are numerical
checks of the proved finite identities, not interval certificates.

The local laboratory exposes every scale, coupling/field controls for the
classical model, selected generated interaction supports on actual dessin
radial lifts, the even-chain $ K,h,C $ trajectory and physical correlation
length, and three fixed quantum matrix calculations. Live quantum parameter
sliders are not offered for uncomputed Hamiltonians. Pair truncation carries
the identity term and reports both partition and probability/density errors.
A late-scale floating-point loss of the small eigenvalue is avoided by
transporting the exact log-eigenratio coordinate $\log r\mapsto 2\log r $;
this is explicitly the transfer RG identity, not a newly diagonalized value.

\subsection{A quantum-feature model trained on actual RG}
The learner targets all 32 unnormalized logweights of the $7\to 5$ spin
$ m=5$ decimation. Its three inputs are $(J_+,J_-,b)$, with ranges
$[0.15,1.1]^2\times[-0.6,0.6]$. Fixed fields are
$ h_j=0.045\cos(1.7 j)+0.02 j/5$. Seed 20261103 fixes 160 input samples, a
120/40 train/test split, and 40 OOD samples with $ J_\pm\in[1.15,1.5]$,
$ b\in[-0.8,0.8]$. The three feature qubits encode the couplings by
$ R_y(1.5 J_+)\otimes R_y(1.5 J_-)\otimes R_y(2 b)|000\rangle $.
Two local $ R_y,R_z $ layers, with a CZ triangle after each, contain 12
nominal optimizer angles. The final three $ R_z $ phases commute with the
measured $ Z $ products and are structurally inactive for this loss; they
still change the displayed state phases. We count them as nominal
parameters, not additional effective expressive directions. Seven nonempty Pauli-$ Z $ products feed an affine
$8\times 32$ readout. There are 268 parameters. Feature registers are not
the seven physical spin sites or five quiver vertices.

The training objective is full-action MSE, including the constant; KL alone
would ignore that constant. A training-only ridge fit initializes the head,
then 120 Adam updates optimize both circuit and head. Actual PennyLane
analytic expectation calls implement the feature circuit. Parameter-shift
derivatives of each feature are combined with the classical head gradient;
full-loss finite differences agree to $6.46\times 10^{-11}$. All data,
parameters, intermediate checkpoints and complete loss curves are preserved.

The held-out full-action MSE is approximately 0.00985, mean
$ D_{\mathrm{KL}}(p_{\rm exact}\Vert p_{\rm learned})=0.00300$, mean TV 0.0290,
and mean absolute log-partition error 0.0610. A classical quadratic-feature
baseline with the same 256 readout parameters obtains 0.00613,0.00240,0.0233
and 0.0138, respectively. OOD action MSE increases to 0.267 for the quantum
hybrid, versus 0.0555 for the baseline. The quantum-feature constant MAE is 0.0444 and higher-body coefficient
RMSE is 0.00900, versus 0.0124 and 0.0130 for the classical baseline; different
metrics favor different aspects. The aggregate results show no quantum
advantage. They expose genuine full-action approximation errors rather
than substituting state-compression success for RG.

Each saved example computes its maximum error over all 32 coarse states;
the proposition then applies to its subsequent exact marginalizations.
These per-example finite bounds are not uniform over the input parameter
box. The synchronized architecture inspector shows literal gate outputs,
Pauli features, head, logweights, generated interactions and partition
value, linked to the exact genus-two dessin/order context. No full order
representation, continuum universality, material-specific accuracy or
arithmetic zero-location bridge is inferred.

\subsection{Quantum TokenGT and the geometric error-correction boundary}
A further executed node/edge-token hybrid uses the same exact RG teacher,
with frozen three-qubit features, trained Q/K/V and a full-action head.
Softmax, componentwise max-plus and hybrid variants are compared with a
matched classical feature model; the latter has lower held-out action error.
An explicit intervention-code/change variational MI objective learns concept
directions; a 61-candidate latent forest and a held-out analogy transfer are
scored on exact physical observables. Rips $ H_1$ is actually reduced over
$\mathbf F_2$, and a continuous Gaussian GFlowNet trains a trajectory-balance
objective with specified root/Lebesgue measures. Empirical residuals do not
prove exact terminal sampling, semantic reasoning or categorical equivalence.

The physical bit-flip repetition channel is separately executed and has
$ p_L=3 p^2-2 p^3$ under ideal recovery. A genuine cellular CSS calculation on
the genus-two one-face dessin gives $[[10,4,1]]$: vertex-incidence rank 6,
zero face boundary, X-distance 1 and Z-distance 4. Thus this cellulation
provides no useful protection merely from its genus. These finite layers,
full formulas, proofs, parameters and scope are developed in the research
book and interactive laboratories; the physical models are not faithful
representations of the infinite completed order and supply no GRH bridge.

\subsection{An actual finite gentle cut for the categorical interface}
\label{rg:cut}
On the genuine $ m=5$ dessin quiver write
$ a_i,b_i:i\to i-1$ with indices modulo 5; cross-branch consecutive pairs
are forbidden. Remove both arrows with source 0. The surviving algebra $ B $
has vertices 0 through 4, arrows $ a_i,b_i $ for $ i=1,\ldots,4$, and relations
$ a_i b_{i-1}=b_i a_{i-1}=0$ for $ i=2,\ldots,4$. It is an explicit finite
gentle quotient, not the completed order itself.

Every surviving positive path stays on one branch and has length at most 4.
For length $\ell=1,\ldots,4$ there are $2(5-\ell)$ paths, while length 0
contributes five idempotents. Therefore $\dim_kB=5+2(4+3+2+1)=25$.
The forbidden-chain resolution already proved for monomial gentle algebras
has two alternating chains of length $ n $ from vertex $ i $ to $ i-n $ when
$1\le n\le i $, and none for $ n>i $. Applying Hom to simples kills its
arrow differentials. Hence $\operatorname{pd}_B S_i=i $ and
$\operatorname{Ext}^4_B(S_4,S_0)\cong k^2$. The split finite-dimensional
simple-module criterion gives $\operatorname{gldim}B=4$.

Every full branch cycle crosses a removed arrow, so the common completed
order parameter $ u $ maps to zero. Thus $ B $ is a finite $ u $-torsion quotient,
not a DVR-free surface order. The difference is visible in derived Hom:
$ S_0$ is projective over $ B $, whereas in the original completed order the
two forbidden arms of length 5 return to vertex 0 and give
$\operatorname{Ext}^5_\Lambda(S_0,S_0)\cong k^2$. The exact restriction along
$\Lambda\twoheadrightarrow B $ is consequently not derived fully faithful.
This supplies a concrete cut on which to investigate graded gentle/Fukaya
models, while proving why its equivalence cannot simply be transported to
the original completed order. Boundary/stops/line-field data and an actual
module-valued neural decoder would still have to be supplied. Exact path
and endpoint enumeration accompanies the prose proof; this cut/category
comparison is not claimed wholly Lean formalized.

